\documentclass[11pt]{article}
\usepackage[margin=1in]{geometry}
\usepackage{amsmath,amssymb,amsthm}
\usepackage{graphicx}
\usepackage[hidelinks]{hyperref}
\setlength{\emergencystretch}{2em}

\newtheorem{theorem}{Theorem}
\newtheorem{lemma}{Lemma}
\newtheorem{corollary}{Corollary}
\newtheorem{remark}{Remark}

\newcommand{\C}{\mathbb C}
\newcommand{\F}{\mathbb F}
\newcommand{\dist}{\operatorname{dist}}
\newcommand{\diag}{\Delta}
\newcommand{\Mcal}{\mathcal M}

\title{Separable Essential Ideals and the UBCP Wall}
\author{Run \texttt{fa\_banach\_001}, lane 19}
\date{June 17, 2026}

\begin{document}
\maketitle

\section*{Source Problem}

The source paper is Jinghao Huang, Karimbergen Kudaybergenov, and Rui Liu,
\emph{The ball-covering property of von Neumann algebras and noncommutative
symmetric spaces}, arXiv:2311.14257.  After proving the commutative unital
criterion for \(C(K)\), it asks for a noncommutative \(C^*\)-algebra version:

\begin{center}
\includegraphics[width=0.92\textwidth]{figures/open_problem_crop.png}
\end{center}

\section*{Definitions}

Let \(X\) be a normed space.  It has BCP if \(S_X\) is covered by countably
many open balls not containing the origin.  Equivalently, there are countably
many centres \(x_n\) such that
\[
        S_X\subseteq\bigcup_n B(x_n,\|x_n\|).
\]
It has UBCP if there is such a cover with bounded radii and with all balls
uniformly separated from a neighbourhood of the origin.

The earlier packet
\[
\texttt{2311.14257\_bcp\_countable\_capture\_exact\_equivalence}
\]
proved that a nonzero normed space \(X\) fails BCP if and only if it has
countable capture: for every countable \(C\subset X\) there is \(u\in S_X\)
such that
\[
        \dist(c,\F u)=\|c\| \qquad(c\in C).
\]

\section*{Main Positive Theorem}

\begin{theorem}[separable essential ideals force UBCP]
\label{thm:essential}
Let \(A\) be a \(C^*\)-algebra.  If \(A\) contains a separable essential
closed two-sided ideal \(I\), then \(A\) has UBCP.
\end{theorem}

\begin{proof}
Since \(I\) is essential, the canonical map from \(A\) into the multiplier
algebra \(\Mcal(I)\) is isometric.  Thus, for \(a\in A\),
\[
        \|a\|
        =
        \sup\{\|ax\|:x\in I,\ \|x\|\le1\}.
\]
Equivalently, if \((e_\alpha)\) is a positive contractive approximate
identity for \(I\), then \(\|ae_\alpha\|\to\|a\|\).  Indeed, choose
\(x\in I\) with \(\|x\|\le1\) and \(\|ax\|\) close to \(\|a\|\), and then use
\(e_\alpha x\to x\).

Let \(D\) be a countable dense subset of the closed unit ball of \(I\), and
put
\[
        D_0=\{d\in D:\|d\|>13/16\}.
\]
For every \(d\in D_0\), take the centre \(2d\) and the fixed radius \(9/8\).
This is a countable family of balls.  Its radii are bounded, and every centre
has norm larger than \(13/8\), so all balls avoid \(B(0,1/2)\).

We show that these balls cover \(S_A\).  Fix \(a\in S_A\).  By the first
paragraph, choose a positive contraction \(e\in I\) such that
\[
        \|ae\|>7/8.
\]
Since \(ae\in I\) and \(\|ae\|\le1\), choose \(d\in D\) with
\[
        \|d-ae\|<1/16.
\]
Then \(\|d\|>7/8-1/16=13/16\), so \(d\in D_0\).

In the unitization of \(\Mcal(I)\), the element \(1-2e\) is a contraction,
because \(e\) is a positive contraction.  Hence
\[
\begin{aligned}
        \|a-2d\|
        &\le \|a-2ae\|+2\|ae-d\|\\
        &=   \|a(1-2e)\|+2\|ae-d\|\\
        &<  1+1/8
        =   9/8.
\end{aligned}
\]
Thus \(a\in B(2d,9/8)\).  The unit sphere is covered by the countable family
above, with uniformly bounded radii and a uniform gap from the origin.
Therefore \(A\) has UBCP.
\end{proof}

\begin{corollary}
\label{cor:no-essential}
If a \(C^*\)-algebra has BCP but not UBCP, then it has no separable essential
ideal.
\end{corollary}

\begin{proof}
This is the contrapositive of Theorem~\ref{thm:essential}, using
UBCP \(\Rightarrow\) BCP.
\end{proof}

\section*{Examples Covered}

The theorem recovers several positive regions by a single argument.
Every separable \(C^*\)-algebra has itself as a separable essential ideal.
If \(H\) is separable, then \(\mathcal K(H)\) is a separable essential ideal
in \(\mathcal B(H)\), so \(\mathcal B(H)\) has UBCP.  If \((A_n)\) is a
sequence of separable \(C^*\)-algebras, then
\[
        \bigoplus_{n=1}^\infty A_n
\]
is a separable essential ideal in \(\prod_n A_n\), so the countable product
has UBCP.  This gives a short common explanation for the countable-reservoir
positive cases in the lane.

\section*{Quantitative UBCP Reformulation}

\begin{lemma}[bounded anti-capture form of UBCP]
\label{lem:anti}
For a normed space \(X\), the following are equivalent:
\begin{enumerate}
    \item \(X\) has UBCP;
    \item there are a bounded countable set \(D\subset X\) and a number
          \(\delta>0\) such that for every \(u\in S_X\) some \(d\in D\)
          satisfies
          \[
                \|u-d\|<\|d\|-\delta .
          \]
\end{enumerate}
Consequently, if \(X\) fails UBCP, then for every bounded countable
\(D\subset X\) and every \(\delta>0\) there is \(u\in S_X\) such that
\[
        \|u-d\|\ge \|d\|-\delta \qquad(d\in D).
\]
\end{lemma}

\begin{proof}
If (2) holds, the balls \(B(d,\|d\|-\delta)\), \(d\in D\), cover \(S_X\).
Their radii are bounded because \(D\) is bounded, and they all avoid
\(B(0,\delta)\).  Hence \(X\) has UBCP.

Conversely, a UBCP cover has bounded radii and avoids some ball
\(B(0,\delta)\).  If its centres are \(d_n\) and its radii are \(r_n\), then
\[
        r_n\le \|d_n\|-\delta
\]
for every \(n\).  Since \(S_X\) is covered and the radii are bounded, the
centres are bounded.  This gives (2), after replacing each \(r_n\) by
\(\|d_n\|-\delta\) if necessary.

The final sentence is the negation of (2).
\end{proof}

\section*{The Remaining Wall: Asymptotic Capture}

\begin{theorem}[BCP but not UBCP gives quotient capture]
\label{thm:quotient}
Let \(X\) be a nonzero normed space over \(\F=\mathbb R\) or \(\mathbb C\).
Assume \(X\) has BCP but not UBCP.  Let \(C=\{c_1,c_2,\ldots\}\subset X\) be
any countable set witnessing the positive side of the exact BCP/capture
equivalence, i.e. for every \(u\in S_X\) there is \(c\in C\) with
\[
        \dist(c,\F u)<\|c\|.
\]
Then the sequence quotient
\[
        X_\infty=\ell_\infty(X)/c_0(X)
\]
contains \(v\in S_{X_\infty}\) such that, for the diagonal embedding
\(\diag:X\to X_\infty\),
\[
        \dist(\diag c,\F v)=\|c\| \qquad(c\in C).
\]
\end{theorem}

\begin{proof}
Let \(Q_\F=\mathbb Q\) in the real case and
\(Q_\F=\mathbb Q+i\mathbb Q\) in the complex case.  For \(m\ge1\), put
\[
        D_m=\{q c_j:1\le j\le m,\ q\in Q_\F,\ |q|\le m\}.
\]
This is a bounded countable subset of \(X\).  Since \(X\) does not have UBCP,
Lemma~\ref{lem:anti} gives \(u_m\in S_X\) such that
\[
        \|u_m-y\|\ge \|y\|-1/m \qquad(y\in D_m).
\]
Let \(v\) be the class of the bounded sequence \((u_m)\) in
\(\ell_\infty(X)/c_0(X)\).  Its quotient norm is
\(\limsup_m\|u_m\|=1\), so \(v\in S_{X_\infty}\).

Fix \(c=c_j\in C\) and \(\lambda\in\F\).  If \(\lambda=0\), then
\(\|\diag c-\lambda v\|=\|c\|\).  Assume \(\lambda\ne0\), and put
\(\mu=\lambda^{-1}\).  Let \(\varepsilon>0\).  Choose \(q\in Q_\F\) with
\[
        |q-\mu|<\varepsilon .
\]
For all sufficiently large \(m\), we have \(m\ge j\) and \(|q|\le m\), hence
\(qc\in D_m\).  Therefore
\[
\begin{aligned}
        \|u_m-\mu c\|
        &\ge \|u_m-qc\|-\|(q-\mu)c\|\\
        &\ge |q|\,\|c\|-1/m-\varepsilon\|c\|\\
        &\ge (|\mu|-2\varepsilon)\|c\|-1/m .
\end{aligned}
\]
Taking the limsup and then letting \(\varepsilon\downarrow0\) gives
\[
        \|v-\mu\diag c\|\ge |\mu|\,\|c\|.
\]
Multiplying by \(|\lambda|\), we obtain
\[
        \|\diag c-\lambda v\|\ge \|c\|.
\]
Since this holds for every scalar \(\lambda\), and equality is obtained at
\(\lambda=0\), the desired distance equality follows.
\end{proof}

\begin{corollary}[shape of a C*-counterexample]
\label{cor:shape}
If a \(C^*\)-algebra \(A\) has BCP but not UBCP, then:
\begin{enumerate}
    \item \(A\) has no separable essential ideal;
    \item for every countable BCP witness \(C\subset A\), the quotient
          \(\ell_\infty(A)/c_0(A)\) contains a unit vector that captures the
          diagonal copy of \(C\).
\end{enumerate}
\end{corollary}

\begin{proof}
The first assertion is Corollary~\ref{cor:no-essential}; the second is
Theorem~\ref{thm:quotient}.
\end{proof}

\section*{What Remains}

Theorem~\ref{thm:essential} is a substantial positive theorem, but it is not
a full solution of Question 3.  It reduces the possible failure of
\[
        \text{BCP}\Rightarrow\text{UBCP}
\]
for \(C^*\)-algebras to algebras with no separable essential ideal.  The most
direct remaining route is therefore:

\begin{center}
\emph{prove that every \(C^*\)-algebra with BCP contains a separable
essential ideal.}
\end{center}

If that statement is true, then Theorem~\ref{thm:essential} immediately gives
the full implication BCP \(\Rightarrow\) UBCP for \(C^*\)-algebras.  If it is
false, Corollary~\ref{cor:shape} says any counterexample must hide the
countable BCP witness in \(A\) while allowing exact line capture of that
witness in the sequence quotient.

\section*{Novelty and Search Bounds}

Local run indexes were searched for arXiv:2311.14257, UBCP, SBCP, essential
ideal, separable essential ideal, multiplier algebra, countable capture, and
the previous lane-19 C*-BCP packets.  The earlier
essential-Fell packet proves a type-I local chart theorem; it does not contain
the general separable-essential-ideal compression proof above.  The
Liu--Liu--Lu--Zheng paper arXiv:2111.04921 was checked for the definitions
and the vector-valued \(C_0(K,X)\) UBCP stability theorem.  No exact version
of Theorem~\ref{thm:essential} or Theorem~\ref{thm:quotient} was found in the
local packet index.

\section*{Human Review Focus}

Reviewers should check:
\begin{enumerate}
    \item the use of essentiality to view \(A\) isometrically inside
          \(\Mcal(I)\);
    \item the passage from arbitrary contractions in \(I\) to positive
          contractions \(e\) using an approximate identity;
    \item the fixed estimate \(\|1-2e\|\le1\);
    \item the quantitative negation of UBCP in Lemma~\ref{lem:anti};
    \item the rational-scalar diagonalization in Theorem~\ref{thm:quotient}.
\end{enumerate}

\begin{thebibliography}{9}

\bibitem{HKL}
Jinghao Huang, Karimbergen Kudaybergenov, and Rui Liu,
\emph{The ball-covering property of von Neumann algebras and noncommutative
symmetric spaces},
arXiv:2311.14257, 2023.

\bibitem{LLLZ}
Minzeng Liu, Rui Liu, Jimeng Lu, and Bentuo Zheng,
\emph{Ball covering property from commutative function spaces to
non-commutative spaces of operators},
arXiv:2111.04921, 2021.

\bibitem{Blackadar}
Bruce Blackadar,
\emph{Operator Algebras: Theory of \(C^*\)-Algebras and Von Neumann
Algebras},
Encyclopaedia of Mathematical Sciences 122, Springer, 2006.

\end{thebibliography}

\end{document}
